\section{Results}
\label{sec:results}

\subsection{Quantitative Results}

\subsubsection{Task Completion}

Figure~\ref{fig:headline} and Table~\ref{tab:text-control-regular} present our headline finding: \textbf{voice agents show substantial drops from text baselines}. Under \textbf{Clean} conditions (studio-quality audio, American accents), the best voice \new{model}\old{provider} already drops 34pp from GPT-5 (51\% vs GPT-5 at 85\%). Under \textbf{Realistic} conditions (background noise, diverse accents, natural turn-taking behaviors), performance drops an additional 12pp to 38\%---voice agents retain only 30--45\% of text SOTA capability. Against non-reasoning text models, the best voice \new{model}\old{provider} nearly matches GPT-4.1 (54\%) under Clean conditions (51\%, just 3pp gap) but still drops 16pp under Realistic.

\begin{table}[h]
\caption{Text vs Voice comparison (pass@1). Text shows GPT-5 (reasoning) and GPT-4.1 (non-reasoning). Voice evaluated under Clean and Realistic conditions. Deltas show gap from GPT-5.}
\label{tab:text-control-regular}
\centering
\begin{small}
\resizebox{\columnwidth}{!}{%
\begin{tabular}{@{}llccc@{}}
\toprule
 &  &  & \multicolumn{2}{c}{\textbf{Voice}} \\
\cmidrule(l){4-5}
\textbf{Domain} & {\textbf{Model}} & \textbf{Text} & \textbf{Clean} & \textbf{Realistic} \\
\midrule
\multirow{3}{*}{All} & {\texttt{gemini-live-2.5}} & \multirow{3}{*}{85\% (54\%)} & 31\% (-54, -63.4\%) & 26\% (-59, -69.5\%) \\
 & {\texttt{gpt-realtime-1.5}} &  & 49\% (-36, -42.1\%) & 35\% (-49, -58.4\%) \\
 & {\texttt{grok-voice}} &  & \textbf{51\% (-34, -40.1\%)} & \textbf{38\% (-46, -54.7\%)} \\
\midrule
\multirow{3}{*}{Retail} & {\texttt{gemini-live-2.5}} & \multirow{3}{*}{81\% (76\%)} & 45\% (-36, -44.8\%) & 30\% (-51, -63.2\%) \\
 & {\texttt{gpt-realtime-1.5}} &  & \textbf{71\% (-10, -12.3\%)} & \textbf{45\% (-36, -44.8\%)} \\
 & {\texttt{grok-voice}} &  & 48\% (-33, -40.4\%) & 39\% (-42, -52.4\%) \\
\midrule
\multirow{3}{*}{Airline} & {\texttt{gemini-live-2.5}} & \multirow{3}{*}{83\% (53\%)} & 28\% (-55, -66.3\%) & 30\% (-53, -63.9\%) \\
 & {\texttt{gpt-realtime-1.5}} &  & \textbf{48\% (-35, -42.2\%)} & \textbf{40\% (-43, -51.8\%)} \\
 & {\texttt{grok-voice}} &  & 46\% (-37, -44.6\%) & 36\% (-47, -56.6\%) \\
\midrule
\multirow{3}{*}{Telecom} & {\texttt{gemini-live-2.5}} & \multirow{3}{*}{90\% (34\%)} & 20\% (-70, -77.6\%) & 18\% (-72, -80.5\%) \\
 & {\texttt{gpt-realtime-1.5}} &  & 28\% (-62, -68.8\%) & 21\% (-69, -76.6\%) \\
 & {\texttt{grok-voice}} &  & \textbf{58\% (-32, -35.7\%)} & \textbf{40\% (-50, -55.2\%)} \\
\bottomrule
\multicolumn{5}{l}{\footnotesize \textit{Text column: GPT-5, reasoning (GPT-4.1, best non-reasoning model). Deltas relative to GPT-5.}} \\
\end{tabular}
}
\end{small}
\end{table}

The Clean-to-Realistic drop varies substantially across \new{models}\old{providers}: just 5pp for \new{Gemini}\old{Google} (8\% of its total voice-text gap) versus 13--14pp for \new{Grok and GPT-realtime}\old{xAI and OpenAI} (roughly one-quarter of their total gap). For most \new{models}\old{providers}, the dominant source of the voice-text gap is the drop from text to Clean voice, not the additional Realistic degradation.


Across \new{models}\old{providers}, \textbf{\new{Grok}\old{xAI} achieves slightly higher scores} (51\% Clean, 38\% Realistic, 2--3pp ahead of \new{GPT-realtime}\old{OpenAI}), while \textbf{\new{Gemini}\old{Google} shows the smallest degradation} under realistic conditions, losing 17\% of its Clean performance compared to 24--28\% for others---a 1.5$\times$ robustness advantage. Domain-specific patterns emerge: \new{Grok}\old{xAI} substantially outperforms others in Telecom (58\% Clean vs 20--28\% for others), while in Retail \new{GPT-realtime}\old{OpenAI} leads with 71\% Clean---the single best per-domain score in the benchmark.

\paragraph{Statistical Reliability.} \new{For all three domains, we conducted 2 independent runs per condition and test statistical significance using paired permutation tests (100k permutations, paired by task ID, Holm-Bonferroni corrected). Pooling across domains (Table~\ref{tab:stat-sig}), the Text $\to$ Realistic and Clean $\to$ Realistic gaps are statistically significant for all three voice models (all $p \le 0.002$), as is the reasoning Text $\to$ Clean gap (all $p < 0.001$). Combined voice models achieve 31--52\% (Clean) and 24--36\% (Realistic), well below text baselines of 56\% (GPT-4.1) and 85\% (GPT-5). The non-reasoning Text $\to$ Clean gap is significant for two of three models; per-domain breakdown and the full pairwise analysis in Appendix~\ref{app:stat-sig}.}\why{addresses beyN W4/Q5}\old{For Retail, where we conducted 2 independent runs per condition, we test statistical significance using paired permutation tests (100k permutations, paired by task ID, Holm-Bonferroni corrected). Both the text-to-Clean gap and the Clean-to-Realistic gap are statistically significant for all three providers (all $p < 0.05$; Table~\ref{tab:stat-sig}). Voice providers achieve 42--68\% (Clean) and 29--43\% (Realistic) across two pooled runs, well below text baselines of 76\% (GPT-4.1) and 82\% (GPT-5). Even the narrowest gap is significant ($p = 0.032$). Full pairwise breakdown in Appendix~\ref{app:stat-sig}.}

\subsubsection{Impact of Acoustic Realism}

To isolate which factors hurt performance most, we conduct ablations on the Retail domain, adding noise, accents, or turn-taking independently (Table~\ref{tab:ablation-single}).

\begin{table}[h]
\caption{Ablation: impact of individual acoustic factors on pass@1 (Retail domain). \new{Each cell shows the absolute pass@1 followed by (absolute pp delta, relative \% delta) from Clean; deltas are computed from un-rounded source values and may not equal subtraction of the displayed integer percentages.}\why{addresses beyN Q1: Table 7 rounding}}
\label{tab:ablation-single}
\centering
\begin{small}
\resizebox{\columnwidth}{!}{%
\begin{tabular}{lccc|c}
\toprule
\textbf{Condition} & {\texttt{gemini-live-2.5}} & {\texttt{gpt-realtime-1.5}} & {\texttt{grok-voice}} & \textbf{All} \\
\midrule
Clean & 45\% & \textbf{71\%} & 48\% & 55\% \\
+ Noise & 40\% (-4, -9.8\%) & \textbf{67\% (-4, -6.2\%)} & 46\% (-2, -3.6\%) & 51\% (-4, -6.4\%) \\
+ Accents & 44\% (-1, -2.0\%) & \textbf{60\% (-11, -16.0\%)} & 30\% (-18, -38.2\%) & 44\% (-10, -18.7\%) \\
+ Turn-taking & 33\% (-11, -25.5\%) & \textbf{57\% (-14, -19.8\%)} & 52\% (+4, +7.3\%) & 47\% (-7, -13.4\%) \\
Realistic & 30\% (-15, -33.3\%) & \textbf{45\% (-26, -37.0\%)} & 39\% (-10, -20.0\%) & 38\% (-17, -31.0\%) \\
\bottomrule
\end{tabular}%
}
\end{small}
\end{table}

\textbf{Accents are the most damaging factor on average}, causing a 10pp average drop, followed closely by turn-taking (7pp) and noise (4pp). However, the accent effect is highly \new{model}\old{provider}-specific: \new{Grok}\old{xAI} is severely affected ($-$18pp, $-$38\% relative), while \new{Gemini}\old{Google} is nearly unaffected ($-$1pp, $-$2\% relative). This finding has accessibility implications, particularly for \new{Grok}\old{xAI} users with non-American accents. Because accents are implemented via TTS personas, these results should be interpreted as indicative rather than definitive.

\textbf{Interactions between factors are complex and \new{model}\old{provider}-specific.} For \new{Gemini}\old{Google}, individual factors sum to $-$16pp and the full Realistic condition causes $-$15pp---a nearly additive interaction where individual effects approximately predict the combined outcome. Turn-taking is \new{Gemini}\old{Google}'s worst single factor ($-$11pp, $-$25\% relative). For \new{Grok}\old{xAI}, the pattern reverses: individual factors sum to $-$16pp, yet the full Realistic drop is only $-$10pp---accents alone devastate \new{Grok}\old{xAI} ($-$18pp), but adding noise and turn-taking partially compensates. Notably, \new{GPT-realtime}\old{OpenAI} has both the highest Clean score (71\%) and the largest relative Realistic degradation ($-$37\%), suggesting that higher baseline capability does not protect against---and may amplify---the impact of speech complexity.

\subsubsection{Voice Interaction Quality}

Beyond task completion, we evaluate conversational dynamics under Realistic conditions (Table~\ref{tab:voice-quality}). We report four aggregate dimensions: \textbf{Latency} (how quickly agents react), \textbf{Responsiveness} (whether agents act when needed), \textbf{Interrupt} (how often agents cut off users mid-speech), and \textbf{Selectivity} (whether agents correctly ignore signals that do not require action).

\begin{table}[h]
\caption{Voice interaction quality (Realistic condition, aggregated across domains). \textbf{Bold} indicates best. Full breakdown in Appendix~\ref{app:voice-metrics-detail}.}
\label{tab:voice-quality}
\centering
\begin{small}
\resizebox{\columnwidth}{!}{%
\begin{tabular}{@{}lcccc@{}}
\toprule
{\textbf{Model}} & \textbf{Latency}$\downarrow$ & \textbf{Responsiveness}$\uparrow$ & \textbf{Interrupt}$\downarrow$ & \textbf{Selectivity}$\uparrow$ \\
\midrule
{\texttt{gemini-live-2.5}} & 1.14s & 69\% & 21\% & 54\% \\
{\texttt{gpt-realtime-1.5}} & \textbf{0.90s} & \textbf{100\%} & \textbf{14\%} & 6\% \\
{\texttt{grok-voice}} & 1.15s & 83\% & 84\% & \textbf{57\%} \\
\bottomrule
\end{tabular}
}
\end{small}
\end{table}



\textbf{\new{GPT-realtime}\old{OpenAI}} excels at latency (0.90s), responsiveness (100\%), and interrupt rate (14\%), but has the worst selectivity (6\%)---responding to nearly all backchannels, vocal tics, and non-directed speech.

\textbf{\new{Grok}\old{xAI}} achieves the best selectivity (57\%), high responsiveness (83\%), and moderate latency (1.15s), but has the highest interrupt rate (84\%)---interrupting users nearly once per turn.

\textbf{\new{Gemini}\old{Google}} has the lowest interrupt rate (21\%), reasonable latency (1.14s), and comparable selectivity to \new{Grok}\old{xAI} (54\%), but the lowest responsiveness (69\%)---failing to respond to nearly a third of user turns.

Each \new{model}\old{provider} excels on a different subset of conversational dimensions but falls short on at least one, revealing a fundamental tradeoff in real-time turn-taking: no current system achieves both reliable responsiveness and appropriate restraint.

\subsection{Qualitative Error Analysis}
\label{sec:analysis}

To characterize failure modes beyond aggregate pass rates---and to verify that observed failures stem from agent behavior rather than artifacts of the benchmark or user simulator---we perform a qualitative error analysis.

\paragraph{Task Selection.} We define $\text{pass}_{\text{text}}$ as tasks where both GPT-4.1 and GPT-5 (reasoning) succeed in text mode, $\text{pass}_{\text{clean}}$ as tasks where a majority of \new{voice models}\old{audio providers} succeed under Clean conditions, and $\text{pass}_{\text{realistic}}$ as tasks where a majority succeed under Realistic conditions. We construct two analysis cohorts:
\begin{itemize}[nosep,leftmargin=*]
    \item \textbf{Voice-Fragile}: Tasks that satisfy $\text{pass}_{\text{text}}$ but not $\text{pass}_{\text{clean}}$, isolating inherent voice interaction challenges.
    \item \textbf{Noise-Fragile}: Tasks that satisfy $\text{pass}_{\text{clean}}$ but not $\text{pass}_{\text{realistic}}$, isolating the impact of acoustic realism (noise, accents, interruptions).
\end{itemize}
For each cohort, we annotate all failed simulations.

\paragraph{Annotation Procedure.} Two independent raters examined each failed simulation, labeling: (1) \textit{error source}---whether the agent or user simulator caused the first critical error; and (2) \textit{error type}---one of logical, transcription, hallucination, VAD/unresponsive, timeout (unresolved in 20 mins), or early termination. Inter-rater agreement was 84\% (76/91 simulations); disagreements were resolved through discussion, reaching 100\% agreement.

\paragraph{Results.} Table~\ref{tab:combined-error-analysis} shows the distribution of error types by source for both cohorts. Full annotations are in Appendix~\ref{app:qualitative-analysis}.

\begin{table}[h]
\caption{Error analysis: distribution of error types by source. Agent errors dominate in both cohorts (79\% and 90\%).}
\label{tab:combined-error-analysis}
\centering
\begin{small}
\resizebox{0.85\columnwidth}{!}{%
\input{results/combined_error_analysis}
}
\end{small}
\end{table}

\textbf{Agent errors dominate}: 79\% of failures in the Voice-Fragile cohort and 90\% in the Noise-Fragile cohort are attributed to the agent rather than the user simulator---suggesting that observed failures primarily reflect agent behavior under our evaluation setup, not simulator artifacts.

\textbf{Logical errors are most common in the Voice-Fragile cohort} (13/43), indicating that voice agents struggle with reasoning even when transcription is accurate. In the \textbf{Noise-Fragile cohort, logical and transcription errors are equally prevalent} (16/48 each), reflecting both reasoning failures and speech recognition errors under noisy conditions.

Authentication is the dominant bottleneck in both cohorts: agents fail to transcribe names and emails even when spelled letter-by-letter, blocking all downstream actions. Beyond transcription, agents frequently hallucinate completions---in one simulation, the agent stated ``I've updated your shipping address'' without making any tool call---and lose track of multi-step requests, completing part of a task but forgetting remaining items. Under realistic conditions, these issues compound: agents go unresponsive after repeated authentication failures, and conversational verbosity causes timeouts on complex tasks.

These qualitative patterns align with the quantitative findings: transcription failures during authentication are consistent with the accent vulnerability observed in the ablations (\S\ref{sec:results}), and the prevalence of hallucinated completions and policy violations even under clean conditions suggests that the voice-text gap is not purely a speech recognition problem---reasoning and grounding challenges persist independently of audio quality.

\new{Re-cutting the 77 agent-attributable failures by the \emph{skill required to succeed}, $\sim$95\% reflect domain-agnostic commonsense conversational skills---spelling, grounding, honesty, multi-part request tracking, and arithmetic---while only $\sim$5\% require domain-specific policy knowledge.}\why{addresses vw95 W2; macs final framing}
